\chapter{Defaults}
\label{ch:defaults}

An important design philosophy of the language is that it has no concept (nor
models) a \emph{null} value or uninitialized memory.

Every variable, constant, function parameter, product types fields etc. have a
\emph{zero}\footnotemark[1] value, which is the value that they are initialized
to at their declaration.

\footnotetext[1]{
	A zero value is also referred to by its more user-friendly name: the
	\emph{default} value; the two terms will be used interchangeably.
}

Every type has a zero value, for example it's the number zero for numeric
types, an empty string for \code{Text}, false for \code{Bool} etc.

\begin{minted}{\giovilexer}
someInteger Int
someInteger -> _out # 0

someText Text
someText -> _out # "" (empty string)

someBool Bool
someBool -> _out # false
\end{minted}

The same applies to non-primitive types: lists and maps (empty list and empty
map respectively), and most importantly to product types, where the zero value
has all its fields initialized to their respective zero values.

\begin{listing}
	\label{lst:product_type_zero_value}
	\caption{
		Example of a variable being initialized with the zero value of its
		product type
	}
	\begin{minted}{\giovilexer}
struct Person <
	name Text
	aliases [Text]
	age Int
	isAdult Bool
	favoriteColor tuple<Int, Int, Int>
>
someDefaultPerson Person
someDefaultPerson.name -> _out # ""
someDefaultPerson.aliaes -> _out # []
someDefaultPerson.age -> _out # 0
someDefaultPerson.isAdult -> _out # false
someDefaultPerson.favoriteColor -> _out # <0,0,0>
	\end{minted}
\end{listing}

For a more formal definition, see~\See{sec:zero_values}.

\clearpage

\section{The default keyword}

The \code{default} keyword can be used in redirections, for example, to "reset"
a variable to its default value

\begin{minted}{\giovilexer}
age Int
69 -> age
age -> _out # 69
default -> age
age -> _out # 0
\end{minted}

You can get the default value of a type with the dot operator

\begin{minted}{\giovilexer}
Int.default -> _out # 0
Bool.default -> _out # false
Person.default -> _out
# Person.<name:"", aliaes:[], age:0, isAdult:false, favoriteColor:<0,0,0>>
\end{minted}

You can also get the default value of a variable, which will be useful later
with \See{sec:default_override}

\begin{minted}{\giovilexer}
age Int
age.default -> _out # 0
\end{minted}

\section{Overriding default}
\label{sec:default_override}

A peculiarity of the language is that the default value is overrideable.

\nb{This is possible only at declaration.}

\begin{minted}{\giovilexer}
age Int default 69
age.default -> _out # 69
\end{minted}

Product types fields can also have their default value overridden

\begin{listing}
\label{lst:product_type_zero_value_overridden}
	\caption{
		Example of a variable being initialized with the overridden zero value
		of its product type
	}\begin{minted}{\giovilexer}
struct Person <
	name Text default "nobody"
	aliases [Text]
	age Int default 18
	isAdult Bool default true
	favoriteColor tuple<Int, Int, Int> default <255, 255, 255>
>
someDefaultPerson Person
someDefaultPerson.name -> _out # nobody
someDefaultPerson.aliaes -> _out # []
someDefaultPerson.age -> _out # 18
someDefaultPerson.isAdult -> _out # true
someDefaultPerson.favoriteColor -> _out # <255,255,255>
	\end{minted}
\end{listing}

